% Nguồn SGK Toán 9 Tập một — Bài 11, trang 66--73:
% https://cdn3.olm.vn/upload/thuvienso/4491668113-page-66-1773022942528.png
% https://cdn3.olm.vn/upload/thuvienso/4491668113-page-67-1773022942853.png
% https://cdn3.olm.vn/upload/thuvienso/4491668113-page-68-1773022943154.png
% https://cdn3.olm.vn/upload/thuvienso/4491668113-page-69-1773022943477.png
% https://cdn3.olm.vn/upload/thuvienso/4491668113-page-70-1773022943773.png
% https://cdn3.olm.vn/upload/thuvienso/4491668113-page-71-1773022944053.png
% https://cdn3.olm.vn/upload/thuvienso/4491668113-page-72-1773022944338.png
% https://cdn3.olm.vn/upload/thuvienso/4491668113-page-73-1773022944721.png
% Nguồn SBT Toán 9 Tập một — Bài 11, PDF trang 41--46:
% SBT TOAN 9 KNTT TAP 1.pdf
% SHA-256: d7ffd04618456682795977e5c47de7874161a35e8606e5bac1906881603a72f7

\section{Tỉ số lượng giác của góc nhọn}

\begin{lesson}
    \subsection{Kiến thức trọng tâm}

    \textbf{Mở đầu}

    Ta có thể xác định ``góc dốc'' $\alpha$ của một đoạn đường dốc khi biết độ dài của dốc là $a$ và độ cao của đỉnh dốc so với đường nằm ngang là $h$ không? (H.4.1). (Trong các tòa chung cư, người ta thường thiết kế đoạn dốc cho người đi xe lăn với góc dốc bé hơn $6^\circ$).

    \begin{center}
    \begin{tikzpicture}
        \coordinate (L) at (0,0);
        \coordinate (R) at (4,0);
        \coordinate (T) at (4,2.2);
        \draw (L)--(R)--(T)--cycle;
        \pic[draw, angle radius=4mm] {right angle=L--R--T};
        \pic[draw, "$\alpha$", angle radius=7mm, angle eccentricity=1.35] {angle=R--L--T};
        \node[above left] at ($(L)!0.5!(T)$) {$a$};
        \node[right] at ($(R)!0.5!(T)$) {$h$};
    \end{tikzpicture}
    \end{center}

    \answerline

    \subsubsection{Khái niệm tỉ số lượng giác của một góc nhọn}

    Cho tam giác $ABC$ vuông tại $A$. Xét góc nhọn $B$. Cạnh $AC$ gọi là cạnh đối của góc $B$, cạnh $AB$ gọi là cạnh kề của góc $B$ (H.4.2).

    \begin{center}
    \begin{tikzpicture}
        \coordinate (B) at (0,0);
        \coordinate (A) at (3.5,0);
        \coordinate (C) at (3.5,2.8);
        \draw (B)--(A)--(C)--cycle;
        \pic[draw, angle radius=4mm] {right angle=B--A--C};
        \node[below left] at (B) {$B$};
        \node[below right] at (A) {$A$};
        \node[above right] at (C) {$C$};
        \node[below] at ($(B)!0.5!(A)$) {cạnh kề của góc $B$};
        \node[right] at ($(A)!0.5!(C)$) {cạnh đối của góc $B$};
        \node[above left] at ($(B)!0.5!(C)$) {cạnh huyền};
    \end{tikzpicture}
    \end{center}

    \textbf{Câu hỏi.} Xét góc $C$ của tam giác $ABC$ vuông tại $A$ (H.4.3). Hãy chỉ ra cạnh đối và cạnh kề của góc $C$.

    \begin{center}
    \begin{tikzpicture}
        \coordinate (A) at (0,0);
        \coordinate (B) at (0,2.8);
        \coordinate (C) at (4,0);
        \draw (A)--(B)--(C)--cycle;
        \pic[draw, angle radius=4mm] {right angle=B--A--C};
        \node[below left] at (A) {$A$};
        \node[above left] at (B) {$B$};
        \node[below right] at (C) {$C$};
    \end{tikzpicture}
    \end{center}

    \answerline

    \textbf{Khái niệm sin, côsin, tang, côtang của góc nhọn $\alpha$}

    \textbf{HĐ1.} Cho tam giác $ABC$ vuông tại $A$ và tam giác $A'B'C'$ vuông tại $A'$ có $\widehat B=\widehat{B'}=\alpha$. Chứng minh rằng:
    \begin{enumerate}[label=\alph*)]
        \item $\triangle ABC\sim\triangle A'B'C'$;
        \item $\dfrac{AC}{BC}=\dfrac{A'C'}{B'C'}$; $\dfrac{AB}{BC}=\dfrac{A'B'}{B'C'}$; $\dfrac{AC}{AB}=\dfrac{A'C'}{A'B'}$; $\dfrac{AB}{AC}=\dfrac{A'B'}{A'C'}$.
    \end{enumerate}

    \answerline
    \answerline

    \textbf{Nhận xét.} Trong Hình 4.4, các tam giác vuông có cùng một góc nhọn $\alpha$ là đồng dạng với nhau. Vì vậy các tỉ số giữa cạnh đối và cạnh huyền (cạnh kề và cạnh huyền), cạnh đối và cạnh kề (cạnh kề và cạnh đối) của góc nhọn $\alpha$ là như nhau, cho dù độ dài các cạnh đối (các cạnh kề) của góc $\alpha$ và các cạnh huyền có thể khác nhau với từng tam giác.

    \begin{center}
    \begin{tikzpicture}
        \coordinate (O) at (0,0);
        \coordinate (P) at (6,0);
        \coordinate (Q) at (6,3.6);
        \coordinate (M) at (3.2,0);
        \path[name path=OQ] (O)--(Q);
        \path[name path=MNline] (M)--++(0,3);
        \path[name intersections={of=OQ and MNline, by=N}];
        \draw (O)--(P);
        \draw (O)--(Q);
        \draw[dashed] (M)--(N);
        \draw[dashed] (P)--(Q);
        \pic[draw, "$\alpha$", angle radius=7mm, angle eccentricity=1.35] {angle=P--O--Q};
        \pic[draw, angle radius=3.5mm] {right angle=O--M--N};
        \pic[draw, angle radius=3.5mm] {right angle=O--P--Q};
    \end{tikzpicture}
    \end{center}

    Cho góc nhọn $\alpha$. Xét tam giác $ABC$ vuông tại $A$ có góc nhọn $B$ bằng $\alpha$ (H.4.5). Ta có:
    \begin{itemize}
        \item Tỉ số giữa cạnh đối và cạnh huyền gọi là sin của $\alpha$, kí hiệu $\sin\alpha$.
        \item Tỉ số giữa cạnh kề và cạnh huyền gọi là côsin của $\alpha$, kí hiệu $\cos\alpha$.
        \item Tỉ số giữa cạnh đối và cạnh kề của góc $\alpha$ gọi là tang của $\alpha$, kí hiệu $\tan\alpha$.
        \item Tỉ số giữa cạnh kề và cạnh đối của góc $\alpha$ gọi là côtang của $\alpha$, kí hiệu $\cot\alpha$.
    \end{itemize}

    \begin{center}
    \begin{tikzpicture}
        \coordinate (B) at (0,0);
        \coordinate (A) at (3.7,0);
        \coordinate (C) at (3.7,3);
        \draw (B)--(A)--(C)--cycle;
        \pic[draw, angle radius=4mm] {right angle=B--A--C};
        \pic[draw, "$\alpha$", angle radius=7mm, angle eccentricity=1.35] {angle=A--B--C};
        \node[below left] at (B) {$B$};
        \node[below right] at (A) {$A$};
        \node[above right] at (C) {$C$};
        \node[below] at ($(B)!0.5!(A)$) {cạnh kề};
        \node[right] at ($(A)!0.5!(C)$) {cạnh đối};
        \node[above left] at ($(B)!0.5!(C)$) {cạnh huyền};
    \end{tikzpicture}
    \end{center}

    \textbf{Chú ý.} Ta có:
    \[
        \sin\alpha=\dfrac{\text{cạnh đối}}{\text{cạnh huyền}},\qquad
        \cos\alpha=\dfrac{\text{cạnh kề}}{\text{cạnh huyền}};
    \]
    \[
        \tan\alpha=\dfrac{\text{cạnh đối}}{\text{cạnh kề}},\qquad
        \cot\alpha=\dfrac{\text{cạnh kề}}{\text{cạnh đối}}.
    \]
    \begin{itemize}
        \item $\cot\alpha=\dfrac{1}{\tan\alpha}$.
        \item $\sin\alpha$, $\cos\alpha$, $\tan\alpha$, $\cot\alpha$ gọi là các tỉ số lượng giác của góc nhọn $\alpha$.
    \end{itemize}

    Sin, côsin của góc nhọn luôn dương và bé hơn $1$ vì trong tam giác vuông, cạnh huyền dài nhất.

    \textbf{Ví dụ 1.} Cho tam giác $ABC$ vuông tại $A$, có $AB=3$ cm, $AC=4$ cm (H.4.6). Hãy tính các tỉ số lượng giác $\sin\alpha$, $\cos\alpha$, $\tan\alpha$ với $\alpha=\widehat B$.

    \begin{center}
    \begin{tikzpicture}
        \coordinate (B) at (0,0);
        \coordinate (A) at (3,0);
        \coordinate (C) at (3,4);
        \draw (B)--(A)--(C)--cycle;
        \pic[draw, angle radius=4mm] {right angle=B--A--C};
        \pic[draw, "$\alpha$", angle radius=7mm, angle eccentricity=1.35] {angle=A--B--C};
        \node[below left] at (B) {$B$};
        \node[below right] at (A) {$A$};
        \node[above right] at (C) {$C$};
        \node[below] at ($(B)!0.5!(A)$) {$3$ cm};
        \node[right] at ($(A)!0.5!(C)$) {$4$ cm};
    \end{tikzpicture}
    \end{center}

    \textbf{Giải}

    Xét $\triangle ABC$ vuông tại $A$, $\widehat B=\alpha$.

    Theo Định lí Pythagore, ta có:
    \[
        BC^2=AC^2+AB^2=4^2+3^2=25
    \]
    nên $BC=5$ (cm).

    Theo định nghĩa của tỉ số lượng giác sin, côsin, tang, ta có
    \[
        \sin\alpha=\dfrac{AC}{BC}=\dfrac{4}{5},\qquad
        \cos\alpha=\dfrac{AB}{BC}=\dfrac{3}{5},\qquad
        \tan\alpha=\dfrac{AC}{AB}=\dfrac{4}{3}.
    \]

    $\sin\alpha$ còn được viết là $\sin\widehat B$ hay $\sin B$. Tương tự cho $\cos\alpha$, $\tan\alpha$ và $\cot\alpha$.

    \textbf{Luyện tập 1.} Cho tam giác $ABC$ vuông tại $A$ có $AB=5$ cm, $AC=12$ cm. Hãy tính các tỉ số lượng giác của góc $B$.

    \answerline
    \answerline

    \textbf{Giá trị lượng giác sin, côsin, tang, côtang của các góc $30^\circ$, $45^\circ$, $60^\circ$}

    \textbf{HĐ2.} Cho tam giác $ABC$ vuông cân tại $A$ và $AB=AC=a$ (H.4.7a).
    \begin{enumerate}[label=\alph*)]
        \item Hãy tính $BC$ và các tỉ số $\dfrac{AB}{BC}$ và $\dfrac{AC}{BC}$. Từ đó suy ra $\sin45^\circ$, $\cos45^\circ$.
        \item Hãy tính các tỉ số $\dfrac{AB}{AC}$ và $\dfrac{AC}{AB}$. Từ đó suy ra $\tan45^\circ$, $\cot45^\circ$.
    \end{enumerate}

    \begin{center}
    \begin{tikzpicture}
        \coordinate (A) at (0,2);
        \coordinate (B) at (-2,0);
        \coordinate (C) at (2,0);
        \draw (A)--(B)--(C)--cycle;
        \pic[draw, angle radius=4mm] {right angle=B--A--C};
        \pic[draw, "$45^\circ$", angle radius=5mm, angle eccentricity=1.35] {angle=A--B--C};
        \pic[draw, "$45^\circ$", angle radius=5mm, angle eccentricity=1.35] {angle=B--C--A};
        \node[above] at (A) {$A$};
        \node[below left] at (B) {$B$};
        \node[below right] at (C) {$C$};
        \node[above left] at ($(A)!0.5!(B)$) {$a$};
        \node[above right] at ($(A)!0.5!(C)$) {$a$};
    \end{tikzpicture}
    \end{center}

    \answerline
    \answerline

    \textbf{HĐ3.} Xét tam giác đều $ABC$ có cạnh bằng $2a$.
    \begin{enumerate}[label=\alph*)]
        \item Tính đường cao $AH$ của tam giác $ABC$ (H.4.7b).
        \item Tính $\sin30^\circ$, $\cos30^\circ$, $\sin60^\circ$ và $\cos60^\circ$.
        \item Tính $\tan30^\circ$, $\cot30^\circ$, $\tan60^\circ$ và $\cot60^\circ$.
    \end{enumerate}

    \begin{center}
    \begin{tikzpicture}
        \coordinate (B) at (-2,0);
        \coordinate (C) at (2,0);
        \coordinate (H) at ($(B)!0.5!(C)$);
        \coordinate (A) at (0,{2*sqrt(3)});
        \draw (A)--(B)--(C)--cycle;
        \draw (A)--(H);
        \pic[draw, angle radius=4mm] {right angle=A--H--C};
        \pic[draw, "$60^\circ$", angle radius=5mm, angle eccentricity=1.35] {angle=A--B--C};
        \pic[draw, "$30^\circ$", angle radius=5mm, angle eccentricity=1.4] {angle=B--A--H};
        \node[above] at (A) {$A$};
        \node[below left] at (B) {$B$};
        \node[below right] at (C) {$C$};
        \node[below] at (H) {$H$};
        \node[below] at ($(B)!0.5!(H)$) {$a$};
        \node[below] at ($(H)!0.5!(C)$) {$a$};
        \node[above right] at ($(A)!0.5!(C)$) {$2a$};
    \end{tikzpicture}
    \end{center}

    \answerline
    \answerline

    Từ HĐ2 và HĐ3, ta có bảng sau:

    \begin{center}
    \begin{tblr}{
        colspec={Q[c] Q[c] Q[c] Q[c]}
    }
        $\alpha$ & $30^\circ$ & $45^\circ$ & $60^\circ$ \\
        $\sin\alpha$ & $\dfrac{1}{2}$ & $\dfrac{\sqrt2}{2}$ & $\dfrac{\sqrt3}{2}$ \\
        $\cos\alpha$ & $\dfrac{\sqrt3}{2}$ & $\dfrac{\sqrt2}{2}$ & $\dfrac{1}{2}$ \\
        $\tan\alpha$ & $\dfrac{\sqrt3}{3}$ & $1$ & $\sqrt3$ \\
        $\cot\alpha$ & $\sqrt3$ & $1$ & $\dfrac{\sqrt3}{3}$ \\
    \end{tblr}
    \end{center}

    \textbf{Ví dụ 2.} Cho tam giác $ABC$ vuông tại $A$ có $\widehat C=30^\circ$ và $BC=a$ (H.4.8). Tính các cạnh $AB$, $AC$ theo $a$.

    \begin{center}
    \begin{tikzpicture}
        \coordinate (A) at (0,0);
        \coordinate (B) at (0,2);
        \coordinate (C) at ({2*sqrt(3)},0);
        \draw (A)--(B)--(C)--cycle;
        \pic[draw, angle radius=4mm] {right angle=B--A--C};
        \pic[draw, "$30^\circ$", angle radius=6mm, angle eccentricity=1.35] {angle=B--C--A};
        \node[below left] at (A) {$A$};
        \node[above left] at (B) {$B$};
        \node[below right] at (C) {$C$};
        \node[above right] at ($(B)!0.5!(C)$) {$a$};
    \end{tikzpicture}
    \end{center}

    \textbf{Giải}

    Ta có $\sin C=\dfrac{AB}{BC}$, suy ra $AB=BC\cdot\sin C=a\cdot\sin30^\circ$.

    Theo bảng trên, $\sin30^\circ=\dfrac{1}{2}$ nên $AB=\dfrac{a}{2}$.

    Tương tự, ta có $\cos C=\dfrac{AC}{BC}$, suy ra $AC=BC\cdot\cos C=a\cdot\cos30^\circ=\dfrac{a\sqrt3}{2}$.

    \textbf{Luyện tập 2.} Cho tam giác $ABC$ vuông tại $A$ có $\widehat C=45^\circ$ và $AB=c$. Tính $BC$ và $AC$ theo $c$.

    \answerline
    \answerline

    \subsubsection{Tỉ số lượng giác của hai góc phụ nhau}

    \textbf{HĐ4.} Cho tam giác $ABC$ vuông tại $C$, có $\widehat A=\alpha$, $\widehat B=\beta$ (H.4.9). Hãy viết các tỉ số lượng giác của góc $\alpha$, $\beta$ theo độ dài các cạnh của tam giác $ABC$. Trong các tỉ số đó, cho biết các cặp tỉ số bằng nhau.

    \begin{center}
    \begin{tikzpicture}
        \coordinate (A) at (-3,0);
        \coordinate (B) at (2,0);
        \coordinate (C) at (0,{sqrt(6)});
        \draw (A)--(B)--(C)--cycle;
        \pic[draw, angle radius=4mm] {right angle=A--C--B};
        \pic[draw, "$\alpha$", angle radius=6mm, angle eccentricity=1.35] {angle=C--A--B};
        \pic[draw, "$\beta$", angle radius=6mm, angle eccentricity=1.35] {angle=A--B--C};
        \node[below left] at (A) {$A$};
        \node[below right] at (B) {$B$};
        \node[above] at (C) {$C$};
    \end{tikzpicture}
    \end{center}

    \answerline
    \answerline

    \begin{keyconcept}
        Nếu hai góc phụ nhau thì sin góc này bằng côsin góc kia, tang góc này bằng côtang góc kia.
    \end{keyconcept}

    \textbf{Chú ý.} Cho $\alpha$ và $\beta$ là hai góc phụ nhau (H.4.9), khi đó:
    \[
        \sin\alpha=\cos\beta,\qquad
        \cos\alpha=\sin\beta,\qquad
        \tan\alpha=\cot\beta,\qquad
        \cot\alpha=\tan\beta.
    \]

    Về số đo, hai góc phụ nhau có thể coi là hai góc nhọn của một tam giác vuông.

    \textbf{Ví dụ 3.} Hãy viết các tỉ số lượng giác sau thành tỉ số lượng giác của góc nhỏ hơn $45^\circ$:
    \[
        \sin60^\circ,\quad \cos75^\circ,\quad \sin52^\circ30',\quad \tan80^\circ,\quad \cot82^\circ.
    \]

    \textbf{Giải}

    Ta có:
    \[
        \sin60^\circ=\cos(90^\circ-60^\circ)=\cos30^\circ;
    \]
    \[
        \cos75^\circ=\sin(90^\circ-75^\circ)=\sin15^\circ;
    \]
    \[
        \sin52^\circ30'=\cos(90^\circ-52^\circ30')=\cos37^\circ30';
    \]
    \[
        \tan80^\circ=\cot(90^\circ-80^\circ)=\cot10^\circ;
    \]
    \[
        \cot82^\circ=\tan(90^\circ-82^\circ)=\tan8^\circ.
    \]

    \textbf{Luyện tập 3.} Hãy giải thích tại sao $\sin35^\circ=\cos55^\circ$, $\tan35^\circ=\cot55^\circ$.

    \answerline

    \subsubsection{Sử dụng máy tính cầm tay tính tỉ số lượng giác của một góc nhọn}

    Về số đo góc, dưới đơn vị độ ($^\circ$) còn có các đơn vị phút ($'$) và giây ($''$) với $1^\circ=60'$, $1'=60''$.

    \textbf{Ví dụ 4.} Dùng MTCT, tính $\sin27^\circ$, $\cos32^\circ15'$, $\tan52^\circ12'$ và $\cot35^\circ23'$ (làm tròn đến chữ số thập phân thứ ba).

    \textbf{Giải}

    \begin{center}
    \begin{tblr}{
        colspec={Q[l] X[l] Q[c]},
        row{1}={font=\bfseries, halign=c}
    }
        Để tính & Bấm phím & Kết quả \\
        $\sin27^\circ$ & \texttt{sin 2 7}\,$^\circ{}'{}''$\,\texttt{=} & $0{,}4539904997$ \\
        $\cos32^\circ15'$ & \texttt{cos 3 2}\,$^\circ{}'{}''$\,\texttt{1 5}\,$^\circ{}'{}''$\,\texttt{=} & $0{,}8457278217$ \\
        $\tan52^\circ12'$ & \texttt{tan 5 2}\,$^\circ{}'{}''$\,\texttt{1 2}\,$^\circ{}'{}''$\,\texttt{=} & $1{,}289192232$ \\
        $\cot35^\circ23'$ & \texttt{tan 3 5}\,$^\circ{}'{}''$\,\texttt{2 3}\,$^\circ{}'{}''$\,\texttt{=} $x^{-1}$ \texttt{=} & $1{,}408003909$ \\
    \end{tblr}
    \end{center}

    Làm tròn đến chữ số thập phân thứ ba ta được $\sin27^\circ\approx0{,}454$; $\cos32^\circ15'\approx0{,}846$; $\tan52^\circ12'\approx1{,}289$; $\cot35^\circ23'\approx1{,}408$.

    \textbf{Lưu ý.} $\cot35^\circ23'=\dfrac{1}{\tan35^\circ23'}$.

    \textbf{Nhận xét.} Để tính $\cot35^\circ23'$, ta có thể tính trực tiếp như trên, hoặc có thể tìm góc phụ với góc $35^\circ23'$ là $54^\circ37'$ rồi dùng MTCT tính $\tan54^\circ37'$ và suy ra kết quả.

    \textbf{Luyện tập 4.} Sử dụng MTCT tính các tỉ số lượng giác và làm tròn kết quả đến chữ số thập phân thứ ba:
    \begin{enumerate}[label=\alph*)]
        \item $\sin40^\circ54'$;
        \item $\cos52^\circ15'$;
        \item $\tan69^\circ36'$;
        \item $\cot25^\circ18'$.
    \end{enumerate}

    \answerline

    \textbf{Chú ý.} Sử dụng MTCT, ta còn có thể tìm được góc khi biết một trong các tỉ số lượng giác của góc đó.

    \textbf{Ví dụ 5.} Dùng MTCT, tìm các góc (làm tròn đến phút) biết $\sin\alpha_1=0{,}3214$, $\cos\alpha_2=0{,}4321$, $\tan\alpha_3=1{,}2742$ và $\cot\alpha_4=1{,}5384$.

    \textbf{Giải}

    \begin{center}
    \begin{tblr}{
        colspec={Q[l] X[l] Q[c] Q[c]},
        row{1}={font=\bfseries, halign=c}
    }
        Biết & Bấm phím & Kết quả & Bấm tiếp \\
        $\sin\alpha_1=0{,}3214$ & \texttt{SHIFT sin 0 . 3 2 1 4 =} & $18{,}74761209$ & $18^\circ44'51{,}4''$ \\
        $\cos\alpha_2=0{,}4321$ & \texttt{SHIFT cos 0 . 4 3 2 1 =} & $64{,}39909458$ & $64^\circ23'56{,}74''$ \\
        $\tan\alpha_3=1{,}2742$ & \texttt{SHIFT tan 1 . 2 7 4 2 =} & $51{,}87495892$ & $51^\circ52'29{,}85''$ \\
        $\cot\alpha_4=1{,}5384$ & \texttt{SHIFT tan 1 . 5 3 8 4} $x^{-1}$ \texttt{=} & $33{,}02491482$ & $33^\circ1'29{,}69''$ \\
    \end{tblr}
    \end{center}

    Làm tròn đến phút ta được $\alpha_1\approx18^\circ45'$; $\alpha_2\approx64^\circ24'$; $\alpha_3\approx51^\circ52'$; $\alpha_4\approx33^\circ1'$.

    \textbf{Chú ý.} Để tìm góc $\alpha$ khi biết $\cot\alpha$, ta có thể tìm góc $(90^\circ-\alpha)$ (vì $\tan(90^\circ-\alpha)=\cot\alpha$) rồi suy ra $\alpha$.

    \textbf{Luyện tập 5.} Dùng MTCT, tìm các góc $\alpha$ (làm tròn đến phút), biết:
    \begin{enumerate}[label=\alph*)]
        \item $\sin\alpha=0{,}3782$;
        \item $\cos\alpha=0{,}6251$;
        \item $\tan\alpha=2{,}154$;
        \item $\cot\alpha=3{,}253$.
    \end{enumerate}

    \answerline

    \textbf{Vận dụng.} Trở lại bài toán ở tình huống mở đầu. Trong một tòa chung cư, biết đoạn dốc vào sảnh tòa nhà dài $4$ m, độ cao của đỉnh dốc bằng $0{,}4$ m.
    \begin{enumerate}[label=\alph*)]
        \item Hãy tính góc dốc.
        \item Hỏi góc đó có đúng tiêu chuẩn của dốc cho người đi xe lăn không?
    \end{enumerate}

    \answerline
    \answerline

    \textbf{Tranh luận.} Để tính khoảng cách giữa hai địa điểm $A$, $B$ không đo trực tiếp được, chẳng hạn $A$ và $B$ là hai địa điểm ở hai bên sông, người ta lấy điểm $C$ về phía bờ sông có chứa $B$ sao cho tam giác $ABC$ vuông tại $B$. Ở bên bờ sông chứa $B$, người ta đo được $\widehat{ACB}=\alpha$ và $BC=a$ (H.4.10). Với các dữ liệu đó, đã tính được khoảng cách $AB$ chưa? Nếu được, hãy tính $AB$, biết $\alpha=55^\circ$, $a=70$ m.

    % TODO(HINH-NGUON): H.4.10 co minh hoa bo song; phan tam giac/toan hoc duoi day duoc dung dung quan he va du kien nguon.
    \begin{center}
    \begin{tikzpicture}
        \coordinate (C) at (0,0);
        \coordinate (B) at (4,0);
        \coordinate (A) at (4,3);
        \draw (C)--(B)--(A)--cycle;
        \pic[draw, angle radius=4mm] {right angle=C--B--A};
        \pic[draw, "$\alpha$", angle radius=7mm, angle eccentricity=1.35] {angle=B--C--A};
        \node[below left] at (C) {$C$};
        \node[below right] at (B) {$B$};
        \node[above right] at (A) {$A$};
        \node[below] at ($(C)!0.5!(B)$) {$a$};
    \end{tikzpicture}
    \end{center}

    \begin{quote}
        Không thể tính được $AB$ vì trong tam giác vuông $ABC$, theo định lí Pythagore, phải biết được hai cạnh mới tính được cạnh thứ ba.
    \end{quote}

    \begin{quote}
        Với các dữ liệu đã biết là có thể tính được khoảng cách $AB$ rồi.
    \end{quote}

    Em hãy cho biết ý kiến của mình.

    \answerline
    \answerline
\end{lesson}


\begin{exercises}
    \subsection{Bài tập}

    \begin{questions}[series=SGK]
        \item Cho tam giác $ABC$ vuông tại $A$. Tính các tỉ số lượng giác sin, côsin, tang, côtang của các góc nhọn $B$ và $C$ khi biết:
        \begin{questions}
            \item $AB=8$ cm, $BC=17$ cm;
            \item $AC=0{,}9$ cm, $AB=1{,}2$ cm.
        \end{questions}

        \answerline
        \answerline

        \item Cho tam giác vuông có một góc nhọn $60^\circ$ và cạnh kề với góc $60^\circ$ bằng $3$ cm. Hãy tính cạnh đối của góc này.

        \answerline

        \item Cho tam giác vuông có một góc nhọn bằng $30^\circ$ và cạnh đối với góc này bằng $5$ cm. Tính độ dài cạnh huyền của tam giác.

        \answerline

        \item Cho hình chữ nhật có chiều dài và chiều rộng lần lượt là $3$ và $\sqrt3$. Tính góc giữa đường chéo và cạnh ngắn hơn của hình chữ nhật (sử dụng bảng giá trị lượng giác trang 69).

        \answerline

        \item
        \begin{questions}
            \item Viết các tỉ số lượng giác sau thành tỉ số lượng giác của các góc nhỏ hơn $45^\circ$:
            \[
                \sin55^\circ,\quad \cos62^\circ,\quad \tan57^\circ,\quad \cot64^\circ.
            \]

            \answerline
            \item Tính $\dfrac{\tan25^\circ}{\cot65^\circ}$, $\tan34^\circ-\cot56^\circ$.

            \answerline
        \end{questions}

        \item Dùng MTCT, tính (làm tròn đến chữ số thập phân thứ ba):
        \begin{questions}
            \item $\sin40^\circ12'$;
            \item $\cos52^\circ54'$;
            \item $\tan63^\circ36'$;
            \item $\cot25^\circ18'$.
        \end{questions}

        \answerline

        \item Dùng MTCT, tìm số đo của góc nhọn $x$ (làm tròn đến phút), biết rằng:
        \begin{questions}
            \item $\sin x=0{,}2368$;
            \item $\cos x=0{,}6224$;
            \item $\tan x=1{,}236$;
            \item $\cot x=2{,}154$.
        \end{questions}

        \answerline
    \end{questions}
\end{exercises}

\begin{practice}
    \subsection{Luyện tập}

    \begin{questions}[series=SBT]
        \item
        \begin{questions}
            \item Vẽ tam giác $ABC$ vuông tại $A$, $AB=3$ cm, $AC=4$ cm. Tính $BC$, $\sin B$, $\cos B$.

            \answerline
            \answerline

            \item Vẽ tam giác $MNP$ vuông tại $M$, $MN=6$ cm, $MP=8$ cm. Hỏi hai tam giác $ABC$, $MNP$ có đồng dạng không? Tính $\sin N$, $\cos N$.

            \answerline
            \answerline
        \end{questions}

        \item
        \begin{questions}
            \item Chứng minh rằng với mọi góc nhọn $\alpha<45^\circ$, ta có
            \[
                \sin(45^\circ-\alpha)=\cos(45^\circ+\alpha),\qquad
                \cos(45^\circ-\alpha)=\sin(45^\circ+\alpha).
            \]

            \answerline
            \answerline

            \item Không dùng MTCT, tính
            \[
                \sin25^\circ+\sin35^\circ+\sin45^\circ-\cos45^\circ-\cos55^\circ-\cos65^\circ.
            \]

            \answerline
        \end{questions}

        \item Khi góc $\alpha$ lần lượt bằng $10^\circ$, $20^\circ$, $30^\circ$, $40^\circ$, hãy dùng MTCT tính $\sin\alpha$ trong mỗi trường hợp (làm tròn đến chữ số thập phân thứ ba).

        \answerline

        \item Hãy dùng MTCT, tìm số đo của góc nhọn $\alpha$ (làm tròn đến độ) trong mỗi trường hợp
        \begin{questions}
            \item Khi $\sin\alpha$ lần lượt bằng $\dfrac{1}{4}$, $\dfrac{1}{3}$, $\dfrac{1}{2}$, $\dfrac{2}{3}$;
            \item Khi $\cos\alpha$ lần lượt bằng $\dfrac{1}{4}$, $\dfrac{1}{3}$, $\dfrac{1}{2}$, $\dfrac{2}{3}$.
        \end{questions}

        \answerline
        \answerline

        \item Biết rằng với mỗi góc nhọn $\alpha$, ta có $\sin^2\alpha+\cos^2\alpha=1$, không dùng MTCT, hãy tính $\sin^225^\circ+\sin^235^\circ+\sin^245^\circ+\sin^255^\circ+\sin^265^\circ$.

        \answerline
        \answerline

        \item Một tam giác vuông có hai cạnh góc vuông đo được $5$ cm, $12$ cm. Hỏi sin của góc nhọn nhỏ nhất của tam giác đó bằng bao nhiêu?

        \answerline

        \item Xét tam giác $ABC$ vuông tại $B$, có $\widehat A=30^\circ$. Tia $Bt$ sao cho $\widehat{CBt}=30^\circ$ cắt tia $AC$ ở $D$, $D$ nằm giữa $A$ và $C$. Chứng minh rằng khoảng cách từ $D$ đến đường thẳng $BC$ bằng $\dfrac{AB}{4}$.

        \answerline
        \answerline

        \item Vẽ góc $\alpha$ trong mỗi trường hợp:
        \begin{questions}
            \item $\cos\alpha=0{,}4$;
            \item $\tan\alpha=\dfrac{2}{3}$;
            \item $\cot\alpha=\dfrac{3}{4}$.
        \end{questions}

        \answerline
        \answerline

        \item Một cái thang dài $3{,}2$ m đặt tựa bức tường, đầu thang đặt đến độ cao $3$ m thì thang tạo với mặt đất góc $\alpha$ xấp xỉ bằng bao nhiêu độ (H.4.6)?

        % TODO(HINH-NGUON): bo sung asset/crop dung Hinh 4.6 SBT trang 45 (cai thang tua tuong, nhan 3,2; 3; alpha); khong redraw xap xi.
        \answerline

        \item Một cái diều có dây diều dài $8$ m, khi dây diều căng thì diều bay ở độ cao $6$ m. Hỏi khi đó dây diều tạo với phương ngang của mặt đất góc nhọn $\alpha$ xấp xỉ bằng bao nhiêu độ (H.4.7)?

        % TODO(HINH-NGUON): bo sung asset/crop dung Hinh 4.7 SBT trang 46 (day dieu 8 m, do cao 6 m, goc alpha); khong redraw xap xi.
        \answerline

        \item Chứng minh tam giác vuông có một góc nhọn có tang bằng $1$ là tam giác vuông cân.

        \answerline
        \answerline

        \item
        \begin{questions}
            \item Tính các góc của tam giác vuông có một góc nhọn có tang bằng $\dfrac{\sqrt3}{3}$.

            \answerline

            \item Một hình chữ nhật có kích thước $3$ và $\sqrt3$. Tính các góc tạo bởi đường chéo và cạnh của hình chữ nhật đó.

            \answerline
        \end{questions}

        \item Dùng MTCT, hãy tìm tang và côtang của góc nhọn $\alpha$ khi $\alpha$ lần lượt bằng $10^\circ$, $20^\circ$, $30^\circ$, $40^\circ$ (làm tròn đến chữ số thập phân thứ ba).

        \answerline

        \item Tính tang, côtang của góc kề đáy của tam giác cân biết cạnh đáy dài $8$ cm, đường cao ứng với đáy dài $5$ cm.

        \answerline
        \answerline

        \item Dùng định nghĩa tỉ số lượng giác $\sin\alpha$, $\cos\alpha$, $\tan\alpha$, $\cot\alpha$ hãy chứng minh rằng:
        \begin{questions}
            \item $\tan\alpha=\dfrac{\sin\alpha}{\cos\alpha}$; $\cot\alpha=\dfrac{\cos\alpha}{\sin\alpha}$;
            \item $1+\tan^2\alpha=\dfrac{1}{\cos^2\alpha}$.
        \end{questions}

        \answerline
        \answerline

        \item Cho góc $\alpha$ có $\tan\alpha=\dfrac{3}{4}$. Tính $\sin\alpha$, $\cos\alpha$.

        \answerline

        \item Với $\alpha<\beta<90^\circ$, chứng minh rằng:
        \begin{questions}
            \item $\cos\alpha>\cos\beta$ (\textit{HD. Sử dụng Ví dụ 5 và bài 4.15});
            \item $\sin\alpha<\sin\beta$ (\textit{HD. Sử dụng công thức $\sin^2\alpha+\cos^2\alpha=1$}).
        \end{questions}

        \answerline
        \answerline
    \end{questions}
\end{practice}

